\documentclass[handout,aspectratio=169]{beamer}

%------------------------------------------------
% PAQUETES
%------------------------------------------------
\usepackage[spanish,es-noshorthands]{babel}
\usepackage[utf8]{inputenc}
\usepackage[T1]{fontenc}
\usepackage{lmodern}
\usepackage{amsmath, amssymb}
\usepackage{tikz}
\usepackage{xcolor}
\usepackage{tcolorbox}
\usepackage{booktabs}
\usepackage{array}
\usepackage{ragged2e}
\usepackage{multicol}

%------------------------------------------------
% TEMA
%------------------------------------------------
\usetheme{Madrid}
\useinnertheme{rounded}

%------------------------------------------------
% COLORES UST
%------------------------------------------------
\definecolor{USTGreen}{HTML}{1B5E20}
\definecolor{USTLightGreen}{HTML}{E8F5E9}
\definecolor{USTDark}{HTML}{263238}
\definecolor{USTGray}{HTML}{F5F5F5}
\definecolor{USTAccent}{HTML}{2E7D32}
\definecolor{AlertRed}{HTML}{C62828}

\setbeamercolor{structure}{fg=USTGreen}
\setbeamercolor{frametitle}{bg=USTGreen,fg=white}
\setbeamercolor{title}{fg=white,bg=USTGreen}
\setbeamercolor{block title}{bg=USTGreen,fg=white}
\setbeamercolor{block body}{bg=USTLightGreen,fg=black}
\setbeamercolor{alerted text}{fg=AlertRed}

\setbeamertemplate{navigation symbols}{}

%------------------------------------------------
% FOOTLINE
%------------------------------------------------
\setbeamertemplate{footline}{
	\leavevmode%
	\hbox{%
		\begin{beamercolorbox}[wd=.82\paperwidth,ht=2.5ex,dp=1ex,left]{author in head/foot}
			\hspace{1em} Fundamentos de Matemática Básica
		\end{beamercolorbox}%
		\begin{beamercolorbox}[wd=.18\paperwidth,ht=2.5ex,dp=1ex,center]{date in head/foot}
			\insertframenumber/\inserttotalframenumber
	\end{beamercolorbox}}%
}

%------------------------------------------------
% CAJAS
%------------------------------------------------
\newtcolorbox{idea}{
	colback=USTLightGreen,
	colframe=USTGreen,
	arc=2mm,
	boxrule=0.8pt,
	title=\textbf{Idea clave}
}

\newtcolorbox{alerta}{
	colback=red!5,
	colframe=AlertRed,
	arc=2mm,
	boxrule=0.8pt,
	title=\textbf{Cuidado}
}

\newtcolorbox{actividad}{
	colback=blue!5,
	colframe=blue!60!black,
	arc=2mm,
	boxrule=0.8pt,
	title=\textbf{Actividad}
}

%------------------------------------------------
% DATOS
%------------------------------------------------
\title[Potencias y Álgebra en los Reales]{Potencias, Expresiones Algebraicas y Patrones}
\subtitle{Nivel universitario}
\author{Prof. Luis Tulio González Alcaino}
\institute{Universidad Santo Tomás \\ Departamento de Ciencias Básicas}
\date{\today}

%------------------------------------------------
\begin{document}
	%------------------------------------------------
	
	%------------------------------------------------
	\begin{frame}
		\titlepage
	\end{frame}
	
	%------------------------------------------------
	\begin{frame}{Ruta de la clase}
		\tableofcontents
	\end{frame}
	
	%------------------------------------------------
	\section{Objetivos de aprendizaje}
	
	\begin{frame}{Objetivos de aprendizaje}
		Al finalizar la clase, el estudiante será capaz de:
		
		\begin{enumerate}
			\item Reconocer el concepto de potencia y sus elementos.
			\item Aplicar correctamente las propiedades de las potencias en los reales.
			\item Identificar términos, expresiones algebraicas y términos semejantes.
			\item Simplificar expresiones algebraicas mediante reducción de términos semejantes.
			\item Integrar patrones y relaciones matemáticas en situaciones algebraicas.
			\item Detectar inconsistencias frecuentes en procedimientos con potencias y álgebra.
		\end{enumerate}
	\end{frame}
	
	%------------------------------------------------
	\section{Concepto de potencia}
	
	\begin{frame}{Concepto de potencia}
		\begin{idea}
			Una potencia representa una multiplicación reiterada de un mismo factor.
		\end{idea}
		
		\[
		a^n = 
		\underbrace{a\cdot a\cdot a\cdots a}_{n\text{ veces}}
		\]
		
		\pause
		
		\begin{block}{Elementos}
			\[
			a^n
			\]
			
			\begin{itemize}
				\item \(a\): base.
				\item \(n\): exponente.
				\item \(a^n\): potencia.
			\end{itemize}
		\end{block}
		
		\pause
		
		\begin{exampleblock}{Ejemplo}
			\[
			3^4 = 3\cdot 3\cdot 3\cdot 3 = 81
			\]
		\end{exampleblock}
	\end{frame}
	
	%------------------------------------------------
	\begin{frame}{Interpretación matemática}
		La potencia permite modelar fenómenos de crecimiento repetitivo, escalamiento y patrones multiplicativos.
		
		\pause
		
		\begin{block}{Ejemplos}
			\[
			2^1=2,\qquad 2^2=4,\qquad 2^3=8,\qquad 2^4=16
			\]
			
			\pause
			
			El patrón es:
			
			\[
			2,\ 4,\ 8,\ 16,\ 32,\ldots
			\]
			
			Cada término se obtiene multiplicando el anterior por \(2\).
		\end{block}
		
		\pause
		
		\begin{actividad}
			Explique qué patrón aparece en la sucesión:
			\[
			3^1,\ 3^2,\ 3^3,\ 3^4,\ 3^5
			\]
		\end{actividad}
	\end{frame}
	
	%------------------------------------------------
	\section{Propiedades de las potencias}
	
	\begin{frame}{Producto de potencias de igual base}
		\begin{idea}
			Cuando se multiplican potencias de igual base, se conserva la base y se suman los exponentes.
		\end{idea}
		
		\[
		a^m\cdot a^n = a^{m+n}
		\]
		
		\pause
		
		\begin{exampleblock}{Ejemplo}
			\[
			x^3\cdot x^5 = x^{3+5}=x^8
			\]
		\end{exampleblock}
		
		\pause
		
		\begin{block}{Justificación}
			\[
			x^3\cdot x^5
			=
			(x\cdot x\cdot x)(x\cdot x\cdot x\cdot x\cdot x)
			=
			x^8
			\]
		\end{block}
	\end{frame}
	
	%------------------------------------------------
	\begin{frame}{Cociente de potencias de igual base}
		\begin{idea}
			Cuando se dividen potencias de igual base, se conserva la base y se restan los exponentes.
		\end{idea}
		
		\[
		\frac{a^m}{a^n}=a^{m-n},
		\qquad a\neq 0
		\]
		
		\pause
		
		\begin{exampleblock}{Ejemplo}
			\[
			\frac{x^7}{x^3}=x^{7-3}=x^4
			\]
		\end{exampleblock}
		
		\pause
		
		\begin{alerta}
			La restricción \(a\neq 0\) es fundamental, porque no existe división por cero.
		\end{alerta}
	\end{frame}
	
	%------------------------------------------------
	\begin{frame}{Potencia de una potencia}
		\begin{idea}
			Cuando una potencia se eleva a otra potencia, se multiplican los exponentes.
		\end{idea}
		
		\[
		(a^m)^n=a^{mn}
		\]
		
		\pause
		
		\begin{exampleblock}{Ejemplo}
			\[
			(x^4)^3=x^{4\cdot 3}=x^{12}
			\]
		\end{exampleblock}
		
		\pause
		
		\begin{block}{Interpretación}
			\[
			(x^4)^3=x^4\cdot x^4\cdot x^4=x^{4+4+4}=x^{12}
			\]
		\end{block}
	\end{frame}
	
	%------------------------------------------------
	\begin{frame}{Potencia de un producto}
		\begin{idea}
			La potencia de un producto se distribuye sobre cada factor.
		\end{idea}
		
		\[
		(ab)^n=a^n b^n
		\]
		
		\pause
		
		\begin{exampleblock}{Ejemplo}
			\[
			(3x)^4=3^4x^4=81x^4
			\]
		\end{exampleblock}
		
		\pause
		
		\begin{alerta}
			No confundir:
			\[
			(3+x)^2 \neq 3^2+x^2
			\]
		\end{alerta}
	\end{frame}
	
	%------------------------------------------------
	\begin{frame}{Exponente cero y exponente negativo}
		\begin{block}{Exponente cero}
			\[
			a^0=1,
			\qquad a\neq 0
			\]
		\end{block}
		
		\pause
		
		\begin{exampleblock}{Ejemplo}
			\[
			7^0=1,\qquad x^0=1\quad (x\neq 0)
			\]
		\end{exampleblock}
		
		\pause
		
		\begin{block}{Exponente negativo}
			\[
			a^{-n}=\frac{1}{a^n},
			\qquad a\neq 0
			\]
		\end{block}
		
		\pause
		
		\begin{exampleblock}{Ejemplo}
			\[
			x^{-3}=\frac{1}{x^3}
			\]
		\end{exampleblock}
	\end{frame}
	
	%------------------------------------------------
	\begin{frame}{Síntesis de propiedades}
		\begin{center}
			\renewcommand{\arraystretch}{1.6}
			\begin{tabular}{>{\centering\arraybackslash}m{4cm} >{\centering\arraybackslash}m{5cm}}
				\toprule
				\textbf{Propiedad} & \textbf{Regla} \\
				\midrule
				Producto igual base & \(a^m a^n=a^{m+n}\) \\
				Cociente igual base & \(\dfrac{a^m}{a^n}=a^{m-n}\) \\
				Potencia de potencia & \((a^m)^n=a^{mn}\) \\
				Potencia de producto & \((ab)^n=a^n b^n\) \\
				Exponente cero & \(a^0=1\) \\
				Exponente negativo & \(a^{-n}=\dfrac{1}{a^n}\) \\
				\bottomrule
			\end{tabular}
		\end{center}
	\end{frame}
	
	%------------------------------------------------
	\section{Errores frecuentes con potencias}
	
	\begin{frame}{Detección de inconsistencias}
		Analicemos los siguientes procedimientos.
		
		\pause
		
		\begin{alerta}
			\[
			x^3+x^5=x^8
			\]
		\end{alerta}
		
		\pause
		
		Esto es incorrecto.
		
		\[
		x^3+x^5
		\]
		
		no puede simplificarse sumando exponentes, porque la propiedad
		
		\[
		x^m\cdot x^n=x^{m+n}
		\]
		
		solo se aplica en multiplicaciones.
		
		\pause
		
		\begin{block}{Conclusión}
			\[
			x^3+x^5
			\]
			
			queda expresado como suma, salvo que exista factorización:
			
			\[
			x^3+x^5=x^3(1+x^2)
			\]
		\end{block}
	\end{frame}
	
	%------------------------------------------------
	\begin{frame}{Más inconsistencias}
		Determine si cada afirmación es verdadera o falsa.
		
		\begin{enumerate}
			\item \((x^2)^5=x^{10}\)
			\item \(x^4\cdot x^3=x^{12}\)
			\item \(\dfrac{x^9}{x^4}=x^5\)
			\item \((2x)^3=2x^3\)
			\item \(x^{-2}=\dfrac{1}{x^2}\)
		\end{enumerate}
		
		\pause
		
		\begin{block}{Respuestas}
			\[
			\begin{array}{c|c}
				\text{Afirmación} & \text{Valor lógico} \\
				\hline
				1 & \text{Verdadera} \\
				2 & \text{Falsa} \\
				3 & \text{Verdadera} \\
				4 & \text{Falsa} \\
				5 & \text{Verdadera}
			\end{array}
			\]
		\end{block}
	\end{frame}
	
	%------------------------------------------------
	\section{Elementos básicos del álgebra}
	
	\begin{frame}{Término algebraico}
		\begin{idea}
			Un término algebraico es una expresión formada por números y letras unidos mediante multiplicación y potencias.
		\end{idea}
		
		\pause
		
		\[
		-5x^3y^2
		\]
		
		\pause
		
		\begin{block}{Elementos}
			\begin{itemize}
				\item Coeficiente numérico: \(-5\).
				\item Parte literal: \(x^3y^2\).
				\item Variables: \(x\), \(y\).
				\item Exponentes: \(3\), \(2\).
			\end{itemize}
		\end{block}
	\end{frame}
	
	%------------------------------------------------
	\begin{frame}{Expresión algebraica}
		\begin{idea}
			Una expresión algebraica es una combinación de términos algebraicos mediante operaciones.
		\end{idea}
		
		\[
		4x^2-7x+3
		\]
		
		\pause
		
		\begin{block}{Ejemplos}
			\[
			3x+5
			\]
			
			\[
			2a^2b-4ab^2+7
			\]
			
			\[
			\frac{x^2-1}{x+3}
			\]
		\end{block}
		
		\pause
		
		\begin{alerta}
			Una expresión algebraica no siempre es una ecuación.  
			Una ecuación debe contener una igualdad.
		\end{alerta}
	\end{frame}
	
	%------------------------------------------------
	\begin{frame}{Términos semejantes}
		\begin{idea}
			Dos términos son semejantes si tienen exactamente la misma parte literal.
		\end{idea}
		
		\pause
		
		\[
		7x^2y
		\qquad\text{y}\qquad
		-3x^2y
		\]
		
		son términos semejantes.
		
		\pause
		
		\[
		5xy^2
		\qquad\text{y}\qquad
		5x^2y
		\]
		
		no son términos semejantes.
		
		\pause
		
		\begin{block}{Criterio}
			Para que sean semejantes:
			
			\begin{itemize}
				\item deben tener las mismas variables;
				\item cada variable debe tener el mismo exponente;
				\item el coeficiente puede ser distinto.
			\end{itemize}
		\end{block}
	\end{frame}
	
	%------------------------------------------------
	\begin{frame}{Reducción de términos semejantes}
		\begin{exampleblock}{Ejemplo}
			Reducir:
			
			\[
			8x^2-3x+5x^2+9x-4
			\]
		\end{exampleblock}
		
		\pause
		
		Agrupamos términos semejantes:
		
		\[
		(8x^2+5x^2)+(-3x+9x)-4
		\]
		
		\pause
		
		Sumamos coeficientes:
		
		\[
		13x^2+6x-4
		\]
		
		\pause
		
		\begin{block}{Resultado}
			\[
			8x^2-3x+5x^2+9x-4
			=
			13x^2+6x-4
			\]
		\end{block}
	\end{frame}
	
	%------------------------------------------------
	\section{Patrones y relaciones algebraicas}
	
	\begin{frame}{Patrones numéricos y algebraicos}
		Observe la siguiente tabla:
		
		\[
		\begin{array}{c|c}
			n & 3n+2 \\
			\hline
			1 & 5 \\
			2 & 8 \\
			3 & 11 \\
			4 & 14 \\
			5 & 17
		\end{array}
		\]
		
		\pause
		
		\begin{idea}
			El patrón aumenta de \(3\) en \(3\), por lo tanto puede modelarse mediante una expresión lineal.
		\end{idea}
		
		\pause
		
		\[
		a_n=3n+2
		\]
		
		\pause
		
		\begin{block}{Interpretación}
			La variable \(n\) representa la posición del término dentro de la secuencia.
		\end{block}
	\end{frame}
	
	%------------------------------------------------
	\begin{frame}{Patrón cuadrático}
		Considere la sucesión:
		
		\[
		1,\ 4,\ 9,\ 16,\ 25,\ 36,\ldots
		\]
		
		\pause
		
		Esta sucesión corresponde a:
		
		\[
		1^2,\ 2^2,\ 3^2,\ 4^2,\ 5^2,\ 6^2,\ldots
		\]
		
		\pause
		
		\begin{block}{Modelo}
			\[
			a_n=n^2
			\]
		\end{block}
		
		\pause
		
		\begin{actividad}
			Determine el término general de la sucesión:
			
			\[
			2,\ 8,\ 18,\ 32,\ 50,\ldots
			\]
		\end{actividad}
		
		\pause
		
		\[
		2=2\cdot1^2,\quad 8=2\cdot2^2,\quad 18=2\cdot3^2
		\]
		
		\[
		a_n=2n^2
		\]
	\end{frame}
	
	%------------------------------------------------
	\begin{frame}{Patrones con potencias}
		Considere:
		
		\[
		5,\ 10,\ 20,\ 40,\ 80,\ldots
		\]
		
		\pause
		
		Cada término se obtiene multiplicando por \(2\).
		
		\pause
		
		\[
		5,\quad 5\cdot 2,\quad 5\cdot 2^2,\quad 5\cdot 2^3,\quad 5\cdot 2^4
		\]
		
		\pause
		
		\begin{block}{Término general}
			Si \(n\) comienza en \(1\), entonces:
			
			\[
			a_n=5\cdot 2^{n-1}
			\]
		\end{block}
		
		\pause
		
		\begin{alerta}
			El exponente es \(n-1\), no \(n\), porque el primer término corresponde a \(2^0\).
		\end{alerta}
	\end{frame}
	
	%------------------------------------------------
	\section{Aplicaciones de nivel universitario}
	
	\begin{frame}{Ejemplo aplicado 1: crecimiento multiplicativo}
		Una población bacteriana en una muestra se duplica cada hora.  
		Inicialmente hay \(P_0=1200\) bacterias.
		
		\pause
		
		\begin{block}{Modelo}
			Después de \(t\) horas:
			
			\[
			P(t)=1200\cdot 2^t
			\]
		\end{block}
		
		\pause
		
		\begin{exampleblock}{Pregunta}
			¿Cuántas bacterias habrá después de \(6\) horas?
		\end{exampleblock}
		
		\pause
		
		\[
		P(6)=1200\cdot 2^6
		\]
		
		\[
		P(6)=1200\cdot 64
		\]
		
		\[
		P(6)=76800
		\]
		
		\pause
		
		\begin{block}{Respuesta}
			Después de \(6\) horas habrá \(76800\) bacterias.
		\end{block}
	\end{frame}
	
	%------------------------------------------------
	\begin{frame}{Ejemplo aplicado 2: simplificación algebraica}
		Simplifique:
		
		\[
		A=4x^3y^2\cdot 3x^2y^5
		\]
		
		\pause
		
		Multiplicamos coeficientes:
		
		\[
		4\cdot 3=12
		\]
		
		\pause
		
		Aplicamos producto de potencias de igual base:
		
		\[
		x^3\cdot x^2=x^5
		\]
		
		\[
		y^2\cdot y^5=y^7
		\]
		
		\pause
		
		\begin{block}{Resultado}
			\[
			A=12x^5y^7
			\]
		\end{block}
	\end{frame}
	
	%------------------------------------------------
	\begin{frame}{Ejemplo aplicado 3: expresión con cocientes}
		Simplifique:
		
		\[
		B=\frac{18x^7y^4}{6x^3y}
		\]
		
		\pause
		
		Reducimos coeficientes:
		
		\[
		\frac{18}{6}=3
		\]
		
		\pause
		
		Restamos exponentes de igual base:
		
		\[
		\frac{x^7}{x^3}=x^4
		\]
		
		\[
		\frac{y^4}{y}=y^3
		\]
		
		\pause
		
		\begin{block}{Resultado}
			\[
			B=3x^4y^3
			\]
		\end{block}
		
		\pause
		
		\begin{alerta}
			Restricción:
			\[
			x\neq 0,\qquad y\neq 0
			\]
		\end{alerta}
	\end{frame}
	
	%------------------------------------------------
	\begin{frame}{Ejemplo aplicado 4: términos semejantes}
		Simplifique:
		
		\[
		C=7a^2b-3ab^2+5a^2b+9ab^2-4b
		\]
		
		\pause
		
		Agrupamos términos semejantes:
		
		\[
		(7a^2b+5a^2b)+(-3ab^2+9ab^2)-4b
		\]
		
		\pause
		
		\[
		12a^2b+6ab^2-4b
		\]
		
		\pause
		
		\begin{block}{Resultado}
			\[
			C=12a^2b+6ab^2-4b
			\]
		\end{block}
	\end{frame}
	
	%------------------------------------------------
	\section{Actividades guiadas}
	
	\begin{frame}{Actividad 1: propiedades de potencias}
		Simplifique las siguientes expresiones:
		
		\[
		\begin{array}{ll}
			\text{a)} & x^4\cdot x^7 \\[0.2cm]
			\text{b)} & \dfrac{a^{10}}{a^3} \\[0.3cm]
			\text{c)} & (y^5)^4 \\[0.2cm]
			\text{d)} & (2x^3y)^2 \\[0.2cm]
			\text{e)} & \dfrac{12x^6y^4}{3x^2y}
		\end{array}
		\]
		
		\pause
		
		\begin{block}{Respuestas}
			\[
			\begin{array}{ll}
				\text{a)} & x^{11} \\[0.2cm]
				\text{b)} & a^7 \\[0.2cm]
				\text{c)} & y^{20} \\[0.2cm]
				\text{d)} & 4x^6y^2 \\[0.2cm]
				\text{e)} & 4x^4y^3
			\end{array}
			\]
		\end{block}
	\end{frame}
	
	%------------------------------------------------
	\begin{frame}{Actividad 2: términos semejantes}
		Reduzca la expresión:
		
		\[
		D=9x^2-5x+4-3x^2+11x-8
		\]
		
		\pause
		
		\[
		D=(9x^2-3x^2)+(-5x+11x)+(4-8)
		\]
		
		\pause
		
		\[
		D=6x^2+6x-4
		\]
		
		\pause
		
		\begin{block}{Resultado}
			\[
			D=6x^2+6x-4
			\]
		\end{block}
	\end{frame}
	
	%------------------------------------------------
	\begin{frame}{Actividad 3: detectar errores}
		Un estudiante realiza el siguiente procedimiento:
		
		\[
		(3x^2)^3=3x^6
		\]
		
		\pause
		
		\begin{alerta}
			El procedimiento es incorrecto.
		\end{alerta}
		
		\pause
		
		La potencia debe aplicarse al coeficiente y a la parte literal:
		
		\[
		(3x^2)^3=3^3(x^2)^3
		\]
		
		\pause
		
		\[
		=27x^6
		\]
		
		\pause
		
		\begin{block}{Resultado correcto}
			\[
			(3x^2)^3=27x^6
			\]
		\end{block}
	\end{frame}
	
	%------------------------------------------------
	\begin{frame}{Actividad 4: análisis de patrón}
		Una secuencia está dada por:
		
		\[
		4,\ 12,\ 36,\ 108,\ 324,\ldots
		\]
		
		\pause
		
		\begin{block}{Pregunta}
			Determine el patrón y proponga una expresión para el término general.
		\end{block}
		
		\pause
		
		Cada término se multiplica por \(3\):
		
		\[
		4,\quad 4\cdot 3,\quad 4\cdot 3^2,\quad 4\cdot 3^3
		\]
		
		\pause
		
		\begin{block}{Término general}
			\[
			a_n=4\cdot 3^{n-1}
			\]
		\end{block}
	\end{frame}
	
	%------------------------------------------------
	\section{Síntesis final}
	
	\begin{frame}{Síntesis conceptual}
		\begin{block}{Potencias}
			Las potencias permiten representar multiplicaciones reiteradas y modelar patrones de crecimiento.
		\end{block}
		
		\pause
		
		\begin{block}{Álgebra en los reales}
			Un término algebraico posee coeficiente, parte literal y exponentes.  
			Una expresión algebraica combina términos mediante operaciones.
		\end{block}
		
		\pause
		
		\begin{block}{Términos semejantes}
			Solo pueden reducirse términos con idéntica parte literal.
		\end{block}
		
		\pause
		
		\begin{block}{Pensamiento matemático}
			Detectar patrones, relaciones e inconsistencias permite justificar procedimientos y evitar errores mecánicos.
		\end{block}
	\end{frame}
	
	%------------------------------------------------
	\begin{frame}{Cierre}
		\begin{center}
			{\Huge \textbf{Conclusión}}
			
			\vspace{0.5cm}
			
			\Large
			Las potencias y el álgebra permiten expresar regularidades, simplificar modelos y analizar relaciones matemáticas de manera rigurosa.
			
			\vspace{0.7cm}
			
			\textbf{Pregunta final:}
			
			\[
			x^2+x^2=2x^2
			\]
			
			pero
			
			\[
			x^2\cdot x^2=x^4
			\]
			
			\pause
			
			¿Por qué ambas operaciones tienen resultados distintos?
		\end{center}
	\end{frame}
	
	%------------------------------------------------
\end{document}